\chapter{Further Reading}
\label{ch:references}

Pointers for going deeper, grouped by topic. Online resources were current at the
time of writing.

\section{The Ada language and its profiles}

\begin{itemize}
\item \textbf{Ada Reference Manual} -- the standard itself.
  \url{http://ada-auth.org/standards/ada22.html}. The most relevant corners for
  this book: \emph{D.13} (the Ravenscar and Jorvik profiles), \emph{D.2.3}
  (priority and dispatching policies), \emph{C.6} (\texttt{Volatile} and
  \texttt{Atomic}), \emph{13.3}/\emph{13.5} (address and record-representation
  clauses), and \emph{7.6} / \texttt{Ada.Finalization} (controlled types).
\item \textbf{John Barnes, \emph{Programming in Ada 2012}} -- the standard
  textbook for the language as a whole.
\item \textbf{AdaCore learn.adacore.com} -- free courses, including
  \emph{Introduction to Embedded Systems Programming} and the
  \emph{Ada for the C++/Java developer} guide. \url{https://learn.adacore.com}.
\end{itemize}

\section{The runtime and the toolchain}

\begin{itemize}
\item \textbf{\texttt{bb-runtimes}} -- AdaCore's bareboard runtime generator that
  this project's runtime is ported from.
  \url{https://github.com/AdaCore/bb-runtimes}. The ESP32-S3 board lives in the
  fork \url{https://github.com/rowsail/bb-runtimes} (branch
  \texttt{xtensa-esp32s3-port}).
\item \textbf{\texttt{svd2ada}} -- generates the typed register packages from an
  SVD. \url{https://github.com/AdaCore/svd2ada}.
\item \textbf{Alire} -- the Ada package manager that supplies the toolchain.
  \url{https://alire.ada.dev}.
\item \textbf{GNAT GPL / GNAT Pro documentation} -- for \texttt{pragma
  Restrictions}, the bareboard runtime structure, and \texttt{gprbuild}.
\end{itemize}

\section{The silicon}

\begin{itemize}
\item \textbf{ESP32-S3 Technical Reference Manual} -- the authoritative register
  reference; every peripheral chapter in this book traces back to it.
  \url{https://www.espressif.com/sites/default/files/documentation/esp32-s3_technical_reference_manual_en.pdf}.
\item \textbf{ESP32-S3 Datasheet} -- pin-out, electrical characteristics, the
  GPIO matrix.
\item \textbf{Espressif SVD files} -- the register descriptions consumed by
  \texttt{svd2ada}. \url{https://github.com/espressif/svd}.
\item \textbf{Xtensa Instruction Set Architecture (LX7)} and the
  \emph{Windowed Register Option} -- background for the context-switch and boot
  chapters (Cadence/Tensilica documentation).
\end{itemize}

\section{Cryptography and TLS}

\begin{itemize}
\item \textbf{SPARKNaCl} -- Roderick Chapman's SPARK\,2014 port of TweetNaCl,
  formally proved free of run-time errors. It supplies the public-key, hashing and
  HMAC/HKDF primitives the pure-Ada TLS client stands on (X25519, the SHA family,
  Ed25519); the entire HTTPS chapter (Chapter~\ref{ch:tls}) rests on it. BSD-3-Clause
  (\textcopyright~2021 Protean Code Limited), vendored under \texttt{crates/sparknacl}.
  \url{https://github.com/rod-chapman/SPARKNaCl}. \emph{Our thanks to Rod Chapman and
  the SPARKNaCl project.}
\item \textbf{NaCl / TweetNaCl} -- the originals SPARKNaCl follows: Daniel~J.
  Bernstein, Tanja Lange, Peter Schwabe and colleagues' compact, misuse-resistant
  cryptography library. \url{https://nacl.cr.yp.to} and
  \url{https://tweetnacl.cr.yp.to}.
\item \textbf{RFC 8446} -- \emph{The Transport Layer Security (TLS) Protocol Version
  1.3}, the specification the TLS chapter implements.
  \url{https://www.rfc-editor.org/rfc/rfc8446}.
\end{itemize}

\section{Conformance and high-integrity Ada}

\begin{itemize}
\item \textbf{ACATS} -- the Ada Conformity Assessment Test Suite.
  \url{http://www.ada-auth.org/acats.html}.
\item \textbf{The Ravenscar Profile for Real-Time Systems} (Burns, Dobbing,
  Vardanega) -- the rationale behind the restricted tasking profiles.
\end{itemize}

\section{This project}

\begin{itemize}
\item \textbf{Source repository} --
  \url{https://github.com/rowsail/ada_esp32s3}. The
  \texttt{QUICKSTART.md} and \texttt{TOOLING.md} there cover setup and the
  developer tooling; the \LaTeX{} source of this book is under \texttt{book/}.
\end{itemize}
